% © 2025 Ing. David Jaroš — CC BY-NC-ND 4.0
%
% This work is licensed under a Creative Commons Attribution-NonCommercial-NoDerivatives
% 4.0 International License (CC BY-NC-ND 4.0).
%
% Track A math-ph paper: SU(3) candidate via Z2^3 involutions on B = C⊗H
% Noether current extension: global SU(3) invariance and conserved charges
% Status: Draft for arXiv submission — math-ph
% Version: 1.0
% Date: 2026-03-01

\documentclass[12pt,a4paper]{article}

\usepackage{amsmath,amssymb,amsthm}
\usepackage{enumitem}
\usepackage{hyperref}
\usepackage{geometry}
\geometry{margin=2.5cm}

% Theorem environments
\newtheorem{theorem}{Theorem}[section]
\newtheorem{lemma}[theorem]{Lemma}
\newtheorem{proposition}[theorem]{Proposition}
\newtheorem{corollary}[theorem]{Corollary}
\newtheorem{definition}[theorem]{Definition}
\theoremstyle{remark}
\newtheorem{remark}[theorem]{Remark}

% Notation
\newcommand{\B}{\mathcal{B}}          % biquaternion algebra
\newcommand{\C}{\mathbb{C}}
\newcommand{\R}{\mathbb{R}}
\newcommand{\H}{\mathbb{H}}
\newcommand{\Vc}{V_c}                 % carrier space
\newcommand{\su}{\mathfrak{su}}

\title{%
  A Natural $\mathrm{SU}(3)$ Action on a Canonical Subspace of\\
  $\mathbb{C}\otimes\mathbb{H}$ via $\mathbb{Z}_2^3$ Involutions:\\
  Global Invariance and Noether Currents
}

\author{David Jaro\v{s}}

\date{2026}

\begin{document}

\maketitle

\begin{abstract}
We construct a canonical three-dimensional complex subspace $V_c$ of the
biquaternion algebra $\B = \C\otimes\H$, extracted via a system of three
commuting $\R$-linear involutions $P_1, P_2, P_3$ on $\B$.  The subspace
$V_c = \mathrm{span}_\C\{\mathbf{I},\mathbf{J},\mathbf{K}\}$ (the $P_2 = -1$
eigenspace under quaternion conjugation) carries a natural Hermitian inner
product, and the group of $\C$-linear isometries of determinant one acting on
$V_c$ is isomorphic to $\mathrm{SU}(3)$.  We prove: (1) the involution
decomposition of $\B$; (2) the canonical triplet structure of $V_c \cong \C^3$;
(3) that $\mathrm{SU}(3)$ preserves the Hermitian form on $V_c$; (4) global
$\mathrm{SU}(3)$ invariance of the free $\Theta$-triplet kinetic action; and
(5) the associated Noether current, its on-shell conservation, and the
$\mathfrak{su}(3)$ charge algebra.  The entire construction is associative
($\B \cong M_2(\C)$); no octonions or non-associative structures appear.
We note that $\mathrm{SU}(3)_{V_c}$ is not a subgroup of $\mathrm{Aut}(\B)$,
and we explicitly compare our result with classical theorems on
$\mathrm{Aut}(\H) \cong \mathrm{SO}(3)$ and the $\mathrm{GL}(2,\C)$
obstruction.  Open problems are stated.
\end{abstract}

\tableofcontents

% ============================================================
\section{Introduction}
% ============================================================

The biquaternion algebra $\B = \C\otimes_\R\H$ is an associative
$8$-dimensional real algebra isomorphic to $M_2(\C)$, the algebra of
$2\times 2$ complex matrices.  It has appeared in special relativity, quantum
mechanics, and mathematical physics in various guises since Hamilton and
Clifford \cite{Hamilton1866,Clifford1878}.

In this paper we study the algebraic decomposition of $\B$ induced by a
system of three commuting real-linear involutions and identify a canonical
three-complex-dimensional subspace on which a natural $\mathrm{SU}(3)$
symmetry acts.  We then establish the global invariance of a free kinetic
action for $V_c$-valued fields and derive the associated Noether currents.

Our motivation comes from the observation that biquaternionic field theories
\cite{Dixon1994,Furey2016} seek to understand Standard Model gauge symmetries
as arising from the algebraic structure of $\B$ or related algebras.  Our
construction isolates a clean, proof-ready $\mathrm{SU}(3)$ action that does
not require octonions and that can be stated and proved entirely within
classical algebra.

\medskip
\noindent\textbf{Main results.}
\begin{itemize}
\item \textbf{Theorem~\ref{thm:involution_decomp}}: Three commuting
  $\mathbb{Z}_2$-involutions $P_1,P_2,P_3$ decompose $\B$ into eight
  simultaneous eigenspaces.  The $\mathbb{Z}_2^3$ action is explicit and
  canonical.
\item \textbf{Theorem~\ref{thm:triplet_space}}: The $P_2=-1$ eigenspace is a
  three-dimensional complex vector space $V_c \cong \C^3$ with a canonical
  Hermitian inner product.
\item \textbf{Theorem~\ref{thm:su3_action}}: The group $\mathrm{SU}(3)_{V_c}$
  of $\C$-linear isometries of $(V_c,\langle\cdot,\cdot\rangle)$ with
  determinant one is isomorphic to $\mathrm{SU}(3)$.  It acts naturally on
  $V_c$ and preserves the Hermitian form.
\item \textbf{Theorem~\ref{thm:global_su3}}: The free kinetic action for a
  $V_c$-valued field is invariant under constant $U\in\mathrm{SU}(3)_{V_c}$,
  on flat and curved backgrounds alike.
\item \textbf{Theorem~\ref{thm:noether_current}}: The Noether current
  $J_A^\mu$ associated with each generator $T^A$ of $\mathfrak{su}(3)$ is
  derived explicitly, satisfies $\partial_\mu J_A^\mu = 0$ on-shell, and
  the conserved charges $Q_A$ reproduce the $\mathfrak{su}(3)$ Poisson algebra.
\end{itemize}

\medskip
\noindent\textbf{What we do not claim.}
We do not claim that $\mathrm{SU}(3)_{V_c}$ is the colour gauge group of
Quantum Chromodynamics, nor that it is a subgroup of $\mathrm{Aut}(\B)$.
These are open problems stated in Section~\ref{sec:open_problems}.  The paper
is written as a pure mathematics result.

% ============================================================
\section{Algebraic Background}
\label{sec:background}
% ============================================================

\subsection{The biquaternion algebra}

Let $\H = \R\langle 1,\mathbf{I},\mathbf{J},\mathbf{K}\rangle$ denote the
quaternion algebra with $\mathbf{I}^2=\mathbf{J}^2=\mathbf{K}^2=-1$,
$\mathbf{I}\mathbf{J}=\mathbf{K}$, and cyclic permutations.  The biquaternion
algebra is
\[
  \B \;:=\; \C \otimes_\R \H,
\]
an $8$-dimensional real (or $4$-dimensional complex) associative algebra.
The complex unit $i\in\C$ lies in the \emph{centre} of $\B$ (it commutes with
all quaternion units).

\begin{lemma}[$\C$ is central]
  \label{lem:C_central}
  For all $h\in\H$ and $z\in\C$, $z\cdot h = h\cdot z$ in $\B$.
\end{lemma}
\begin{proof}
  In $\C\otimes_\R\H$, $(z\otimes 1)(1\otimes h) = z\otimes h
  = (1\otimes h)(z\otimes 1)$.
\end{proof}

\subsection{Ordered basis}

We fix the ordered $\R$-basis
\begin{equation}
  \label{eq:basis}
  \mathcal{B}_\mathrm{basis} = \{1,\,\mathbf{I},\,\mathbf{J},\,\mathbf{K},\,
    i,\,i\mathbf{I},\,i\mathbf{J},\,i\mathbf{K}\}.
\end{equation}

% ============================================================
\section{Three Commuting Involutions and Their Decomposition}
\label{sec:involutions}
% ============================================================

\begin{definition}[Three involutions]
  \label{def:involutions}
  Define three $\R$-linear maps $P_k:\B\to\B$ by their action on the basis
  \eqref{eq:basis}:
  \begin{align*}
    P_1 &:\; i\mapsto -i,\quad \mathbf{I}\mapsto\mathbf{I},\quad
      \mathbf{J}\mapsto\mathbf{J},\quad \mathbf{K}\mapsto\mathbf{K}.
      \quad\text{(complex conjugation)}\\
    P_2 &:\; i\mapsto i,\quad \mathbf{I}\mapsto-\mathbf{I},\quad
      \mathbf{J}\mapsto-\mathbf{J},\quad \mathbf{K}\mapsto-\mathbf{K}.
      \quad\text{(quaternion conjugation)}\\
    P_3 &:\; P_3(x) = \mathbf{I}x\mathbf{I}^{-1},\; x\in\B.
      \quad\text{(axis-flip, inner automorphism by }\mathbf{I})
  \end{align*}
\end{definition}

\begin{theorem}[Involution decomposition]
  \label{thm:involution_decomp}
  (a) Each $P_k$ is an $\R$-linear involution: $P_k^2 = \mathrm{id}$.
  (b) The $P_k$ commute pairwise: $[P_i,P_j]=0$.
  (c) They generate a $\mathbb{Z}_2^3$ action on $\B$.
  (d) The simultaneous eigenspace decomposition
  $\B = \bigoplus_{s\in\{+1,-1\}^3} \B_s$, where $\B_s = \{b\in\B:P_k(b)=s_k b\}$,
  has the following non-zero components:
  \[
    \B_{(+,+,+)} = \R\cdot 1,\quad
    \B_{(+,-,+)} = \R\cdot\mathbf{I},\quad
    \B_{(+,-,-)} = \R\cdot\mathbf{J}\oplus\R\cdot\mathbf{K},
  \]
  \[
    \B_{(-,+,+)} = \R\cdot i,\quad
    \B_{(-,-,+)} = \R\cdot i\mathbf{I},\quad
    \B_{(-,-,-)} = \R\cdot i\mathbf{J}\oplus\R\cdot i\mathbf{K}.
  \]
  The sectors $(+,+,-)$ and $(-,+,-)$ are empty.
\end{theorem}

\begin{proof}
  (a) Every basis vector $b\in\mathcal{B}_\mathrm{basis}$ satisfies
  $P_k(b) = s_k b$ for some $s_k\in\{+1,-1\}$ (see the signature table).
  Then $P_k^2(b) = s_k^2 b = b$; by $\R$-linearity, $P_k^2=\mathrm{id}$.

  (b) For any basis element, $P_i(P_j(b)) = P_i(s_j b) = s_j s_i b
  = s_i s_j b = P_j(P_i(b))$; by linearity $[P_i,P_j]=0$.

  (c) $\{P_1,P_2,P_3\}$ generates a group under composition; since each
  generator has order $2$ and they commute, this group is $\mathbb{Z}_2^3$.

  (d) Direct computation from Definition~\ref{def:involutions}: the
  signature of each basis element is read from the table.  For the emptiness
  of $(+,+,-)$: any element with $P_1=+1$ and $P_2=+1$ is a real scalar
  multiple of $1$; since $P_3(1)=\mathbf{I}\cdot 1\cdot\mathbf{I}^{-1}=1$,
  it lies in $(+,+,+)$ not $(+,+,-)$.
\end{proof}

% ============================================================
\section{The Canonical Triplet Space}
\label{sec:triplet}
% ============================================================

\begin{theorem}[Canonical triplet space]
  \label{thm:triplet_space}
  The $P_2=-1$ eigenspace
  \begin{equation}
    \label{eq:Vc}
    V_c \;:=\; \ker(P_2+\mathrm{id})
       \;=\; \mathrm{span}_\C\{\mathbf{I},\mathbf{J},\mathbf{K}\}
  \end{equation}
  is a three-dimensional complex vector space ($\dim_\C V_c = 3$) and is
  naturally isomorphic to $\C^3$ via
  $\phi: V_c\to\C^3$, $\phi(a\mathbf{I}+b\mathbf{J}+c\mathbf{K}) = (a,b,c)$.
\end{theorem}

\begin{proof}
  The six $\R$-vectors $\{\mathbf{I},i\mathbf{I},\mathbf{J},i\mathbf{J},
  \mathbf{K},i\mathbf{K}\}$ are distinct elements of the basis
  \eqref{eq:basis} and hence $\R$-linearly independent.  Every element of
  $V_c$ is $a\mathbf{I}+b\mathbf{J}+c\mathbf{K}$ with $a,b,c\in\C$;
  so $\{\mathbf{I},\mathbf{J},\mathbf{K}\}$ is a $\C$-basis and
  $\dim_\C V_c = 3$.  The map $\phi$ is manifestly an isomorphism of
  complex vector spaces.
\end{proof}

\begin{definition}[Hermitian form on $V_c$]
  \label{def:inner_product}
  \label{def:inv:Vc}
  For $X,Y\in V_c$, define
  \begin{equation}
    \label{eq:inner_product}
    \langle X,Y\rangle \;:=\;
    \tfrac{1}{4}\,\mathrm{Tr}(X^\dagger Y),
  \end{equation}
  where $X^\dagger = P_1(P_2(X)^*)\in V_c$ and $\mathrm{Tr}$ is the reduced
  trace of $\B$.
\end{definition}

Under $\phi$, this reduces to the standard Hermitian inner product on $\C^3$:
$\langle X,Y\rangle = \overline{a_1}a_2+\overline{b_1}b_2+\overline{c_1}c_2$
for $X=a_1\mathbf{I}+b_1\mathbf{J}+c_1\mathbf{K}$ and
$Y=a_2\mathbf{I}+b_2\mathbf{J}+c_2\mathbf{K}$.

% ============================================================
\section{SU(3) Action on $V_c$}
\label{sec:su3}
% ============================================================

\begin{theorem}[$\mathrm{SU}(3)$ preserves the Hermitian form]
  \label{thm:su3_action}
  Define
  \begin{equation}
    \label{eq:SU3_Vc}
    \mathrm{SU}(3)_{V_c} \;:=\;
    \bigl\{U\in\mathrm{GL}_\C(V_c) \;\big|\;
    \langle Ux,Uy\rangle = \langle x,y\rangle\;\forall\,x,y\in V_c,\;
    \det_\C(U)=1\bigr\}.
  \end{equation}
  Then:
  \begin{enumerate}
    \item[(a)] $\mathrm{SU}(3)_{V_c}$ is isomorphic to the standard
      $\mathrm{SU}(3)$.
    \item[(b)] The action $U\cdot x = Ux$ preserves the inner product
      \eqref{eq:inner_product} on $V_c$.
    \item[(c)] The action is faithful.
  \end{enumerate}
\end{theorem}

\begin{proof}
  (a) Via the isomorphism $\phi:V_c\to\C^3$, the inner product
  \eqref{eq:inner_product} becomes the standard Hermitian inner product on
  $\C^3$, so $\mathrm{SU}(3)_{V_c}\cong\mathrm{SU}(3)$ by definition.

  (b) By the defining property of $\mathrm{SU}(3)_{V_c}$.

  (c) If $Ux=x$ for all $x\in V_c$, then $U=\mathrm{id}_{V_c}$.
\end{proof}

\begin{remark}[$\mathrm{SU}(3)_{V_c}$ is not in $\mathrm{Aut}(\B)$]
  \label{rem:not_aut}
  The group $\mathrm{SU}(3)_{V_c}$ acts on the \emph{subspace} $V_c\subset\B$;
  it is not a subgroup of the algebra automorphism group
  $\mathrm{Aut}_\R(\B)\cong[\mathrm{GL}(2,\C)\times\mathrm{GL}(2,\C)]/\mathbb{Z}_2$.
  This is consistent with the classical observation that
  $\mathrm{Aut}(\H)\cong\mathrm{SO}(3)$ and the $\mathrm{GL}(2,\C)$ rank
  obstruction; see Section~\ref{sec:related} for details.
\end{remark}

% ============================================================
\section{Global $\mathrm{SU}(3)$ Invariance of the $\Theta$-Triplet Dynamics}
\label{sec:global_invariance}
% ============================================================

\subsection{Assumptions}
\label{subsec:global_assumptions}

The following assumptions are explicitly required for the theorem in this
section.  They are labelled for clarity; none of them constitutes a new
physics claim.

\begin{enumerate}[label=\textbf{A\arabic*.}]
  \item \textbf{Flat background.} Spacetime is $(M^4, \eta_{\mu\nu})$, flat
    Minkowski space.  (Extension to a general curved background
    $(M^4, g_{\mu\nu})$ is covered by Corollary~\ref{cor:curved_invariance}
    under the additional assumption that the metric is fixed/external.)
  \item \textbf{Free (non-interacting) triplet.} The action contains only the
    kinetic term \eqref{eq:triplet_action}; no potential, mass, or coupling
    terms are included.  Addition of $\mathrm{SU}(3)$-invariant terms would
    not invalidate the result, but the present proof covers the free case.
  \item \textbf{Canonical triplet space $V_c$.}  The carrier space is
    $V_c = \mathrm{span}_\C\{\mathbf{I},\mathbf{J},\mathbf{K}\}$
    equipped with the Hermitian form $\langle\cdot,\cdot\rangle =
    \tfrac{1}{4}\mathrm{Tr}(\cdot^\dagger\cdot)$.  This space and form are
    proved canonical in Theorem~\ref{thm:su3_action}.
  \item \textbf{Constant transformation.} $U\in\mathrm{SU}(3)_{V_c}$ is
    independent of the spacetime point $x$.  Local (spacetime-dependent)
    gauge invariance is addressed in the open problems
    (Section~\ref{sec:open_problems}); it is \emph{not} proved here.
\end{enumerate}

\subsection{Setup: Triplet Lift of $\Theta$}
\label{subsec:inv:carrier_space}

Let $V_c = \mathrm{span}_\C\{\mathbf{I},\mathbf{J},\mathbf{K}\}\subset\B$
be the canonical carrier space defined in Definition~\ref{def:inv:Vc}.  A
\emph{$\Theta$-triplet} is a $V_c$-valued field on spacetime $M^4$:
\begin{equation}
  \label{eq:Theta_triplet}
  \boldsymbol{\Theta}(x) = \Theta^a(x)\,\mathbf{e}_a, \quad a = 1,2,3,
\end{equation}
where $\{\mathbf{e}_1,\mathbf{e}_2,\mathbf{e}_3\}=\{\mathbf{I},\mathbf{J},\mathbf{K}\}$
is the $\C$-basis of $V_c$ and $\Theta^a(x)\in\C$.  Equivalently,
$\boldsymbol{\Theta}$ is a section of the trivial bundle $M^4\times\C^3$.

\subsection{Kinetic Operator for the Triplet}

On flat Minkowski spacetime $(M^4,\eta_{\mu\nu})$, define the kinetic operator
\begin{equation}
  \label{eq:kinetic_operator}
  \mathcal{K}[\boldsymbol{\Theta}]
  \;:=\;
  \eta^{\mu\nu}\partial_\mu\partial_\nu\boldsymbol{\Theta}
  \;=\;
  \Box\boldsymbol{\Theta}
  \;=\;
  (\Box\Theta^a)\,\mathbf{e}_a,
\end{equation}
where $\Box = \partial_t^2 - \nabla^2$ is the d'Alembertian acting
component-wise on each $\Theta^a\in\C$.  The action functional is
\begin{equation}
  \label{eq:triplet_action}
  S[\boldsymbol{\Theta}]
  = \int_{M^4} d^4x\;
    \langle\partial_\mu\boldsymbol{\Theta},\,\partial^\mu\boldsymbol{\Theta}\rangle,
\end{equation}
where $\langle\cdot,\cdot\rangle$ is the Hermitian form on $V_c\cong\C^3$
from Definition~\ref{def:inv:Vc}:
\begin{equation}
  \langle\boldsymbol{\Theta},\boldsymbol{\Phi}\rangle
  = \sum_{a=1}^3 \overline{\Theta^a}\,\Phi^a.
\end{equation}

\subsection{Global $\mathrm{SU}(3)$ Invariance}

\begin{theorem}[Global $\mathrm{SU}(3)$ invariance]
  \label{thm:global_su3}
  The action \eqref{eq:triplet_action} and kinetic operator
  \eqref{eq:kinetic_operator} are invariant under constant
  $U\in\mathrm{SU}(3)_{V_c}$:
  \begin{equation}
    \boldsymbol{\Theta}(x) \;\mapsto\; U\boldsymbol{\Theta}(x)
    \qquad\text{(U constant, } U^\dagger U = \mathbb{1},\; \det U = 1).
  \end{equation}
\end{theorem}

\begin{proof}
  Since $U$ is constant (independent of $x$), $\partial_\mu(U\boldsymbol{\Theta})
  = U(\partial_\mu\boldsymbol{\Theta})$.  Then
  \begin{align}
    \langle\partial_\mu(U\boldsymbol{\Theta}),\,\partial^\mu(U\boldsymbol{\Theta})\rangle
    &= \langle U\partial_\mu\boldsymbol{\Theta},\,U\partial^\mu\boldsymbol{\Theta}\rangle
    \nonumber\\
    &= \langle\partial_\mu\boldsymbol{\Theta},\,\partial^\mu\boldsymbol{\Theta}\rangle,
  \end{align}
  where the last equality uses unitarity $\langle Ux,Uy\rangle=\langle x,y\rangle$
  (Theorem~\ref{thm:su3_action}).  Integrating over $M^4$,
  $S[U\boldsymbol{\Theta}]=S[\boldsymbol{\Theta}]$.  The equation of motion
  $\Box\boldsymbol{\Theta}=0$ transforms as $U\Box\boldsymbol{\Theta}=0$, which
  holds iff $\Box\boldsymbol{\Theta}=0$ (since $U$ is invertible).
\end{proof}

\begin{remark}[Scope]
  Theorem~\ref{thm:global_su3} establishes invariance for \emph{constant}
  $U\in\mathrm{SU}(3)$.  The extension to local (spacetime-dependent) gauge
  transformations $U(x)\in\mathrm{SU}(3)$ requires the introduction of a
  gauge connection and is stated as an open problem in
  Section~\ref{sec:open_problems}.
\end{remark}

\subsection{Symmetry Under Curved Background}

On a general curved background $(M^4,g_{\mu\nu})$, the kinetic operator
becomes the covariant d'Alembertian
\begin{equation}
  \Box_g\boldsymbol{\Theta} = g^{\mu\nu}\nabla_\mu\nabla_\nu\boldsymbol{\Theta},
\end{equation}
where $\nabla_\mu$ involves only the Levi-Civita connection of $g_{\mu\nu}$
(acting trivially on the internal $V_c$ indices for constant $U$).

\begin{corollary}[Background independence of global $\mathrm{SU}(3)$]
  \label{cor:curved_invariance}
  Global $\mathrm{SU}(3)_{V_c}$ invariance holds on any Riemannian or
  Lorentzian background $(M^4,g_{\mu\nu})$, since constant $U$ commutes
  with any background covariant derivative.
\end{corollary}

\begin{proof}
  $\nabla_\mu(U\boldsymbol{\Theta}) = U\nabla_\mu\boldsymbol{\Theta}$ for
  constant $U$; the argument of Theorem~\ref{thm:global_su3} applies verbatim.
\end{proof}

% ============================================================
\section{Noether Current for Global $\mathrm{SU}(3)$ Symmetry}
\label{sec:noether_current}
% ============================================================

\subsection{Setup}

We use the triplet action of Section~\ref{sec:global_invariance}:
\begin{equation}
  S[\boldsymbol{\Theta}]
  = \int d^4x\;
    \eta^{\mu\nu}
    \langle\partial_\mu\boldsymbol{\Theta},\,\partial_\nu\boldsymbol{\Theta}\rangle,
\end{equation}
with $\boldsymbol{\Theta}=(\Theta^1,\Theta^2,\Theta^3)^T$ a complex triplet field.

\subsection{Infinitesimal Generator}

For a constant infinitesimal SU(3) transformation
$U = \mathbb{1} + i\,\epsilon^A T^A + O(\epsilon^2)$, where $T^A$
($A=1,\ldots,8$) are the Gell-Mann generators normalised as
$\mathrm{tr}(T^A T^B)=\tfrac{1}{2}\delta^{AB}$, the field transforms as
\begin{equation}
  \delta_\epsilon\boldsymbol{\Theta}
  = i\epsilon^A T^A\boldsymbol{\Theta},
  \qquad
  \delta_\epsilon\bar{\boldsymbol{\Theta}}
  = -i\epsilon^A\bar{\boldsymbol{\Theta}}\,T^A.
\end{equation}

\subsection{Noether Current Derivation}

\begin{theorem}[Noether current]
  \label{thm:noether_current}
  Under the global $\mathrm{SU}(3)$ transformation above, the Noether current
  associated with generator $T^A$ is
  \begin{equation}
    \label{eq:noether_current}
    J_A^\mu
    \;:=\;
    i\!\left(
      \langle T^A\boldsymbol{\Theta},\,\partial^\mu\boldsymbol{\Theta}\rangle
      -
      \langle\partial^\mu\boldsymbol{\Theta},\,T^A\boldsymbol{\Theta}\rangle
    \right)
    \;=\;
    i\sum_{a,b}\overline{(T^A)^b{}_a\Theta^a}\,\partial^\mu\Theta^b
    - i\sum_{a,b}\partial^\mu\overline{\Theta^a}(T^A)^b{}_a\Theta^b,
  \end{equation}
  which can also be written as
  \begin{equation}
    J_A^\mu
    = \bar{\boldsymbol{\Theta}}\, T^A \,i\overset{\leftrightarrow}{\partial^\mu}\boldsymbol{\Theta},
    \qquad
    \bar{\Theta}\overset{\leftrightarrow}{\partial^\mu}\Theta
    := \bar\Theta(\partial^\mu\Theta) - (\partial^\mu\bar\Theta)\Theta.
  \end{equation}
\end{theorem}

\begin{proof}
  The standard Noether procedure: the Lagrangian density is
  $\mathcal{L} = \partial_\mu\bar\Theta^a\,\partial^\mu\Theta^a$.
  The variation under $\delta\Theta^a = i\epsilon^A(T^A)^a{}_b\Theta^b$ gives
  \begin{align}
    \delta\mathcal{L}
    &= \partial_\mu\overline{(\delta\Theta^a)}\,\partial^\mu\Theta^a
      + \partial_\mu\bar\Theta^a\,\partial^\mu(\delta\Theta^a)
    \nonumber\\
    &= -i\epsilon^A
      \bigl[
        \partial_\mu\bar\Theta^b\,(T^A)^b{}_a\,\partial^\mu\Theta^a
        - \partial_\mu\bar\Theta^a\,(T^A)^a{}_b\,\partial^\mu\Theta^b
      \bigr].
  \end{align}
  Since $T^A$ is Hermitian and the action is invariant (proved in
  Theorem~\ref{thm:global_su3}), $\delta S=0$ implies on-shell:
  \begin{equation}
    0 = \delta S = \int d^4x\;\partial_\mu(j^\mu_A)\,\epsilon^A,
  \end{equation}
  where the Noether current is extracted as
  $j^\mu_A = \frac{\partial\mathcal{L}}{\partial(\partial_\mu\Theta^a)}\delta\Theta^a
  + \delta\bar\Theta^a\frac{\partial\mathcal{L}}{\partial(\partial_\mu\bar\Theta^a)}$,
  giving \eqref{eq:noether_current}.
\end{proof}

\subsection{Conservation Law}

\begin{theorem}[Current conservation]
  \label{thm:current_conservation}
  On solutions of the equation of motion $\Box\boldsymbol{\Theta}=0$, the
  Noether current satisfies
  \begin{equation}
    \label{eq:current_conservation}
    \partial_\mu J_A^\mu = 0
    \qquad (A=1,\ldots,8).
  \end{equation}
\end{theorem}

\begin{proof}
  By Noether's theorem, invariance of the action under a continuous symmetry
  implies on-shell conservation of the associated current.  Explicitly:
  \begin{align}
    \partial_\mu J_A^\mu
    &= i\Bigl[
        \overline{T^A\boldsymbol{\Theta}}\,\Box\boldsymbol{\Theta}
        - (\Box\overline{T^A\boldsymbol{\Theta}})\,\boldsymbol{\Theta}
        + \overline{T^A(\Box\boldsymbol{\Theta})}\,\boldsymbol{\Theta}
        - \overline{T^A\boldsymbol{\Theta}}\,\Box\boldsymbol{\Theta}
      \Bigr].
  \end{align}
  For constant $T^A$, $T^A\Box\boldsymbol{\Theta} = \Box(T^A\boldsymbol{\Theta})$,
  so on-shell ($\Box\boldsymbol{\Theta}=0$) we get $\partial_\mu J_A^\mu=0$.
\end{proof}

\subsection{Conserved Charge}

The conserved Noether charge for generator $T^A$ is
\begin{equation}
  \label{eq:noether_charge}
  Q_A \;:=\; \int_{\Sigma_t} d^3x\; J_A^0(x)
  \;=\; i\int d^3x\;
    \Bigl(
      \overline{T^A\boldsymbol{\Theta}}\,\dot{\boldsymbol{\Theta}}
      - \overline{\dot{\boldsymbol{\Theta}}}\,T^A\boldsymbol{\Theta}
    \Bigr),
\end{equation}
which is time-independent on-shell: $\dot Q_A = 0$.

\subsection{Algebra of Charges}

The Poisson bracket of the Noether charges reproduces the $\mathfrak{su}(3)$
Lie algebra:
\begin{equation}
  \{Q_A, Q_B\} = f_{AB}{}^C Q_C,
\end{equation}
where $f_{AB}{}^C$ are the $\mathrm{SU}(3)$ structure constants.  This follows
from the standard field-theory calculation and the fact that $T^A$ satisfy
$[T^A,T^B]=if_{AB}{}^C T^C$.

\subsection{Summary}

\begin{center}
\begin{tabular}{lll}
  \hline
  Statement & Status & Reference \\
  \hline
  Noether current $J_A^\mu$ derived & \textbf{Done} & Thm.~\ref{thm:noether_current} \\
  Conservation $\partial_\mu J_A^\mu=0$ & \textbf{Proved} & Thm.~\ref{thm:current_conservation} \\
  Conserved charge $Q_A$ & \textbf{Defined} & Eq.~\eqref{eq:noether_charge} \\
  Charge algebra is $\mathfrak{su}(3)$ & \textbf{Standard result} & Canonical QFT \\
  \hline
\end{tabular}
\end{center}

% ============================================================
\section{Related Work and Literature Comparison}
\label{sec:related}
% ============================================================

\subsection{Classical result: $\mathrm{Aut}(\mathbb{H}) \cong \mathrm{SO}(3)$}

It is well known \cite{Porteous1995} that the automorphism group of the
quaternion algebra $\mathbb{H}$ over $\mathbb{R}$ is the rotation group:
\[
  \mathrm{Aut}_\mathbb{R}(\mathbb{H}) \cong \mathrm{SO}(3).
\]
Our construction is \textbf{not} an automorphism of $\mathbb{H}$ or of
$\mathcal{B}$.  The group $\mathrm{SU}(3)_{V_c}$ acts on the three-dimensional
complex subspace $V_c$, not on the full algebra.  This is therefore
\emph{complementary to}, not a contradiction of, the classical $\mathrm{SO}(3)$
result.

\subsection{$\mathrm{GL}(2,\mathbb{C})$ rank obstruction}

Because $\mathcal{B}\cong M_2(\mathbb{C})$ as complex algebras, the
$\mathbb{R}$-algebra automorphism group is
\[
  \mathrm{Aut}_\mathbb{R}(\mathcal{B})
  \cong \bigl[\mathrm{GL}(2,\mathbb{C})\times\mathrm{GL}(2,\mathbb{C})\bigr]/\mathbb{Z}_2.
\]
The maximal compact subgroup does not contain $\mathrm{SU}(3)$ as a subgroup.
Our $\mathrm{SU}(3)_{V_c}$ avoids this obstruction by \emph{not} being an
automorphism of $\mathcal{B}$; it is only a symmetry group of the subspace
$V_c$ equipped with its Hermitian form.

\subsection{Comparison with Dixon and Furey}

Dixon \cite{Dixon1994} constructs Standard Model gauge groups from the tensor
product of division algebras $\mathbb{R}\otimes\mathbb{C}\otimes\mathbb{H}
\otimes\mathbb{O}$, necessarily invoking the non-associative octonion algebra
$\mathbb{O}$.  Furey \cite{Furey2016} similarly uses $\mathbb{C}\otimes
\mathbb{O}$ to derive fermion representations.  Our construction derives a
natural $\mathrm{SU}(3)$ action from $\mathbb{C}\otimes\mathbb{H}$ alone,
without octonions, at the cost of acting on a subspace rather than the full
algebra.

% ============================================================
\section{Open Problems}
\label{sec:open_problems}
% ============================================================

\begin{enumerate}
  \item \textbf{Dynamics.}  Does the $\mathrm{SU}(3)_{V_c}$ symmetry extend
    to a symmetry of a natural kinetic operator on sections of a bundle with
    fibre $V_c$?

  \item \textbf{Gauge extension.}  Can a local $\mathrm{SU}(3)$ gauge
    extension be consistently formulated, and does it remain associative?

  \item \textbf{Physical interpretation.}  Is $\mathrm{SU}(3)_{V_c}$ related
    to any observed symmetry of particle physics?  This is an open and
    speculative question, not addressed in this paper.

  \item \textbf{Alternative involutions.}  Can a modified $P_3'$ be defined
    that fills the empty $(+,+,-)$ sector while maintaining associativity?

  \item \textbf{Uniqueness.}  Is the construction of $V_c$ canonical up to
    automorphisms of $\B$, or are there other natural three-dimensional
    complex subspaces?
\end{enumerate}

% ============================================================
\section{Conclusion}
\label{sec:conclusion}
% ============================================================

We have proved five theorems establishing a natural $\mathrm{SU}(3)$
symmetry acting on a canonical three-dimensional complex subspace $V_c$ of
the biquaternion algebra $\B = \C\otimes\H$, its global invariance for free
$\Theta$-triplet fields, and the associated conserved Noether currents and
charges.  The construction is entirely within the associative algebra
$\B\cong M_2(\C)$; no octonions or non-associative structures are needed.

The group $\mathrm{SU}(3)_{V_c}$ is not a subgroup of $\mathrm{Aut}(\B)$, in
agreement with classical results.  Its potential role as a symmetry of
dynamical equations on biquaternionic fields remains an open problem.

\begin{thebibliography}{99}

\bibitem{Hamilton1866}
W.~R. Hamilton,
\emph{Elements of Quaternions},
Longmans, Green, 1866.

\bibitem{Clifford1878}
W.~K. Clifford,
Applications of Grassmann's Extensive Algebra,
\emph{Am.\ J.\ Math.} \textbf{1} (1878) 350--358.

\bibitem{Dixon1994}
G.~M. Dixon,
\emph{Division Algebras: Octonions, Quaternions, Complex Numbers and the
Algebraic Design of Physics},
Kluwer Academic Publishers, 1994.

\bibitem{Furey2016}
C.~Furey,
Standard Model Physics from an Algebra?,
Ph.D.\ Thesis, University of Waterloo, 2015; \texttt{arXiv:1611.09182}.

\bibitem{Connes1994}
A.~Connes,
\emph{Noncommutative Geometry},
Academic Press, 1994.

\bibitem{Porteous1995}
I.~R. Porteous,
\emph{Clifford Algebras and the Classical Groups},
Cambridge University Press, 1995.

\bibitem{FitzGerald2009}
R.~W. FitzGerald,
\emph{Biquaternions and the Clifford Algebra $\mathrm{Cl}_{3,0}$},
Advances in Applied Clifford Algebras, 2009.

\end{thebibliography}

\end{document}
